\section{Mergeable Summaries as a Local-Law Subclass}
\label{app:proof-mergeable}

This appendix makes the topology in the mergeable-summary comparison explicit.
A classical mergeable-summary guarantee is usually stated with an implicit merge
tree: encode the shards at the leaves, merge states up an allowed binary tree,
and query the root. C-TreePO writes the tree \(T\) explicitly because the local
laws are checked on the realized leaf and merge calls. The clean containment is
therefore a same-tree statement first, followed by the stronger all-schedules
statement.

\subsection{Same-Tree Embedding}

Fix a leaf partition \(D_1,\ldots,D_n\) and a binary tree \(T\) over those
leaves. A classical state-level summary supplies
\[
  \mathrm{build}:\mathrm{Stream}(\alpha)\to S,\qquad
  \mathrm{merge}:S\times S\to S,\qquad
  \mathrm{query}:S\to Y,
\]
together with a validity relation \(\mathrm{Valid}(D,s)\). The state-level
obligations are:
\[
  \mathrm{Valid}(D,\mathrm{build}(D)),
  \qquad
  \mathrm{Valid}(D_1,s_1)\wedge \mathrm{Valid}(D_2,s_2)
  \Rightarrow
  \mathrm{Valid}(D_1\Vert D_2,\mathrm{merge}(s_1,s_2)),
\]
and correctness of the root score for valid states.

We first keep these obligations in exactly this Agarwal form, before
translating them into C-TreePO language. The raw validity and mergeability
obligations are:
\begin{center}
\small
\begin{tabular}{l}
\path|agarwalOriginal_validity_of_mergeTree|\\
\path|agarwalOriginal_readout_of_mergeTree|\\
\path|agarwalOriginal_build_valid_sized_state|\\
\path|agarwalOriginal_merge_valid_sized_state|\\
\path|agarwalOriginal_mergeTree_valid_sized_state|
\end{tabular}
\end{center}
These are just the paper transition in compact notation: valid leaf summaries
and valid-state mergeability imply a valid root state for every binary merge
tree, and valid-state query correctness makes the root query answer the oracle
on the concatenated stream. The sized variants are the fixed-\(\varepsilon\)
statements for \(S(D,\varepsilon)\) and \(k(|D|,\varepsilon)\).

Run this scheme on the same tree \(T\). Leaves store
\(s_v=\mathrm{build}(D_v)\). Internal nodes store
\(s_v=\mathrm{merge}(s_L,s_R)\). Then C1 is exactly leaf validity/score
preservation, and C3 is exactly validity closure under the state merge at the
corresponding internal node. If the C-Tree state map is the serialized
sketch-state operation, this is the same computation with different notation.

C2 is the extra C-TreePO reuse law. It is automatic only when the sketch/codec
also has a stable re-entry witness: passing a produced state through the
summary interface again does not change the task-relevant score. This is why
the bridge uses explicit sketch-local witnesses, not just associativity of the
state merge. In the appendix notation, the same-tree inclusion is:
\[
\texttt{sameTree\_sketch\_to\_local\_laws}
\]
with the reusable statement
\[
\texttt{local\_laws\_bundle\_of\_sketch}.
\]
The exact and approximate routes are
\[
\texttt{sketchCodec\_nests\_exactLocalLaws},\qquad
\texttt{sketchCodecApprox\_nests\_approxLocalLaws}.
\]

\subsection{Validity Recovered by Exact and Approximate Laws}

The C-TreePO reverse direction is a validity statement, not a state-identity
statement. The validity relation is:
\[
\texttt{CTreeSummaryValid}\ f^\star\ x\ z
\quad\equiv\quad
D(f^\star(z),f^\star(x))=0,
\]
with the epsilon variant
\[
\texttt{CTreeSummaryValidWithin}\ f^\star\ \varepsilon\ x\ z
\quad\equiv\quad
D(f^\star(z),f^\star(x))\le \varepsilon.
\]
For distributions, the supportwise and expected forms are:
\begin{center}
\small
\begin{tabular}{l}
\path|CTreeSummaryPMFValid|\\
\path|CTreeSummaryPMFValidWithin|\\
\path|CTreeSummaryPMFExpectedValidWithin|
\end{tabular}
\end{center}

The Agarwal side is represented by the corresponding valid-summary sets:
\[
\texttt{AgarwalValidSummarySet}
\qquad\text{and}\qquad
\texttt{AgarwalValidSizedSummarySet}.
\]
The first is the set of states satisfying the chosen
\(\mathrm{Valid}_{\varepsilon}(D,s)\) relation. The second also records the
paper's \(k(|D|,\varepsilon)\) size bound.

Exact laws give supportwise validity. C1 and C3 on a fixed tree imply that the
one-pass root distribution is supported on valid summaries for the represented
input:
\[
\texttt{exactLocalLaws\_root\_validSummaryPMF}.
\]
With C2 included, the same statement holds for multi-round \(Z_R\) reductions:
\[
\texttt{exactLocalLaws\_multiround\_validSummaryPMF}.
\]
C2 is also the re-entry statement for produced summaries:
\[
\texttt{exactLocalLaws\_resummary\_valid}.
\]
This is the same claim in the notation above: if the local laws hold exactly,
the root summary is valid.

If an external bridge maps C-Tree oracle-validity into the selected Agarwal
validity relation, the same conclusion becomes set containment:
\begin{center}
\small
\begin{tabular}{l}
\path|exactLocalLaws_root_support_subset_agarwalValidSummarySet|\\
\path|exactLocalLaws_multiround_support_subset_agarwalValidSummarySet|\\
\path|exactLocalLaws_root_support_subset_agarwalValidSizedSummarySet|\\
\path|exactLocalLaws_multiround_support_subset_agarwalValidSizedSummarySet|
\end{tabular}
\end{center}
The sized versions use the supplied Agarwal size-profile hypothesis; they do
not derive \(k(|D|,\varepsilon)\) from local laws.

Combining this with the classical valid-state query guarantee gives the exact
query equivalence used in Proposition~\ref{prop:exact-query-equivalence}.
If the C-Tree \(T\) and the classical merge tree \(t\) represent the same input,
then every realized C-Tree output has the same oracle value as the classical
root query:
\[
\texttt{exactLocalLaws\_stateLevelReadout\_equivalence}.
\]
The conclusion is deliberately equality of oracle/query values, not
identity of internal states.

The set-level version quantifies over every Agarwal-valid state for the same
represented stream:
\begin{center}
\small
\begin{tabular}{l}
\path|exactLocalLaws_agarwalValidSummarySet_readout_equivalence|\\
\path|exactLocalLaws_agarwalValidSizedSummarySet_readout_equivalence|\\
\path|exactLocalLaws_stateLevelValidSummarySet_readout_equivalence|
\end{tabular}
\end{center}
Thus exact equivalence is not tied to a single chosen merge tree: any state in
the valid \(S(D,\varepsilon)\) set has the same root score.

Approximate laws have two distinct readings. If one assumes a supportwise
epsilon law, supportwise epsilon-validity follows:
\begin{center}
\small
\begin{tabular}{l}
\path|supportwiseApproxLocalLaws_root_validSummaryPMF|\\
\path|supportwiseApproxLocalLaws_multiround_validSummaryPMF|
\end{tabular}
\end{center}
The audited aggregate bundle used in the paper is weaker: it bounds expected
root distortion through the composed C1/C3/C2 budget. The corresponding statement
is:
\[
\texttt{approxLocalLaws\_expectedValidSummaryPMF},
\]
and, when the bundle is certified below a target \(\varepsilon\),
\[
\texttt{approxLocalLaws\_certifiedAtEpsilon\_expectedValidSummaryPMF}.
\]
Thus aggregate audits certify expected or probabilistic error statements; they
do not by themselves say that every randomized support point is
\(\varepsilon\)-valid.

Finally, Agarwal's \(S(D,\varepsilon)\) notation also includes a size profile
\(k(|D|,\varepsilon)\). Local laws recover oracle/query validity. A
\(\texttt{ValidSizedState}\) follows only after supplying an external bridge
from C-Tree validity to the chosen state-validity relation and an external size
profile. The sized variants make this explicit:
\begin{center}
\small
\begin{tabular}{l}
\path|exactLocalLaws_root_validSizedState_of_external_sizeProfile|\\
\path|exactLocalLaws_multiround_validSizedState_of_external_sizeProfile|
\end{tabular}
\end{center}

\subsection{State-Level Nesting}

The Agarwal-style state-level statement does not require scalar query answers
to merge. The state carries whatever information is needed for the root query;
the query is applied after all state merges have finished. This is the crucial
distinction from a strict oracle-homomorphism statement.

The relational phrasing records that distinction. A state-level summary
with build validity, merge validity, and valid-state query correctness
instantiates C-TreePO's C1/C3 local-law fragment and its relational shape.
At the local-law level, the exact C1 and C3 facts are:
\[
\texttt{agarwalOriginal\_C1\_localLaw},\qquad
\texttt{agarwalOriginal\_C3\_localLaw}.
\]
Together they recover the root decision statement
\[
\texttt{agarwalOriginal\_localLaw\_readout\_of\_mergeTree}.
\]
At the relational C-TreePO level, the same hypotheses give:
\begin{center}
\small
\begin{tabular}{l}
\path|agarwalOriginal_nests_relational_ctreepo|\\
\path|stateLevelSummary_nests_relational_ctreepo|
\end{tabular}
\end{center}
For any merge tree \(t\), the root score then follows as
\begin{center}
\small
\begin{tabular}{l}
\path|agarwalOriginal_readout_of_mergeTree|\\
\path|stateLevelSummary_root_readout_nests_ctreepo|
\end{tabular}
\end{center}
Metric and randomized variants keep their own guarantees:
\begin{center}
\small
\begin{tabular}{l}
\path|agarwalOriginal_readout_error_of_mergeTree|\\
\path|agarwalOriginal_nests_epsilon_relational_ctreepo|\\
\path|epsilonStateLevelSummary_root_readout_nests_ctreepo|\\
\path|randomizedStateLevelSummary_nests_relational_ctreepo|
\end{tabular}
\end{center}
Thus an \(\varepsilon\)-valid summary remains \(\varepsilon\)-valid
for the query, and a randomized summary with root-validity probability
at least \(p\) gives C-TreePO query success with the same probability. No exact
deterministic equality is asserted in those cases.

\subsection{Local-Law Subspace Inclusion}

Let \(\mathrm{MS}(f)\) be the class of operators induced by sketch/codec
systems carrying the exact leaf, merge, and compatibility witnesses above.
For a fixed tree \(T\), let \(\mathrm{ExactLL}(T,f)\) be the exact C1/C2/C3
local-law subspace. The inclusion is
\[
  \mathrm{MS}(f) \subseteq \mathrm{ExactLL}(T,f),
\]
recorded as
\[
\texttt{mergeableSketch\_nests\_exactLocalLawSubspace}.
\]
Inside a chosen neural-operator class \(\mathcal C\), the overlap
\(\mathcal C\cap\mathrm{MS}(f)\) is contained in the exact local-law neural
operator subspace:
\[
\texttt{mergeableSketchOverlap\_nests\_exactLocalLawNeuralOperators}.
\]
This is a membership claim, not an optimization guarantee. A neural operator
does not become mergeable merely by being expressive; it must either supply
these witnesses or be certified by the local-law audit.

\subsection{Fixed-\texorpdfstring{\(\varepsilon\)}{epsilon} NO/FNO Recovery}

The fixed-target learned route is stated as expected-validity. The condition
\[
\texttt{EpsilonGoodLearnedSummary}
\]
means that the \(Z_R\) root distribution has expected oracle distortion at most
the selected target \(\varepsilon\).

For ordinary neural-operator approximation, the expected-validity statement is:
\[
\texttt{neuralOperator\_fixedEpsilon\_expectedValidSummaryPMF}.
\]
Its hypotheses are exactly the ingredients used in
Proposition~\ref{prop:fixed-epsilon-no-recovery}: an exact ideal operator in
the class, a realized neural operator in the class, a compact realized-call
approximation bridge, and a transferred local-law budget below the fixed target.
The conclusion is expected-\(\varepsilon\) validity for the learned
multi-round tree.

For the FNO route, the finite-dimensionalization argument supplies the same
kind of realized-call bridge:
\[
\texttt{fno\_fixedEpsilon\_expectedValidSummaryPMF}.
\]
Use this route when the approximation witness comes from the
finite-mode/Fourier neural-operator certificate rather than a generic
uniform neural-operator approximation assumption.

When local laws are checked through a learned score \(\widehat f\), the true
oracle statement includes the two-sided proxy-oracle gap slack:
\begin{center}
\small
\begin{tabular}{l}
\path|calibratedNeuralOperator_fixedEpsilon_expectedValidSummaryPMF|\\
\path|calibratedFNO_fixedEpsilon_expectedValidSummaryPMF|
\end{tabular}
\end{center}
These statements require the total proxy-oracle gap budget, namely the local-law
budget plus \(2\varepsilon_{\mathrm{orc}}\), to be at most the final target
\(\varepsilon\). They still certify expected validity, not supportwise validity
of every randomized realization.

The same idea can be phrased existentially as "learn an
\(\varepsilon\)-good version":
\begin{center}
\small
\begin{tabular}{l}
\path|exists_neuralOperator_fixedEpsilon_goodLearnedSummaryPMF|\\
\path|exists_fno_fixedEpsilon_goodLearnedSummaryPMF|\\
\path|exists_calibratedNeuralOperator_fixedEpsilon_goodLearnedSummaryPMF|\\
\path|exists_calibratedFNO_fixedEpsilon_goodLearnedSummaryPMF|
\end{tabular}
\end{center}
These variants do not say that gradient descent finds the operator. They say
that any approximation or learning guarantee that returns a realized NO/FNO with
the certified budget immediately yields an \(\varepsilon\)-good learned C-Tree
summary.

The proxy-oracle gap assumption is given its legacy Lean name:
\[
\texttt{OracleRecoveredWithin}.
\]
It is a name for uniform recovery of \(\fstar\) by the learned readout
\(\widehat f\). Under this assumption, the true-oracle budget reads:
\begin{center}
\small
\begin{tabular}{l}
\path|neuralOperator_oracleRecovery_expectedValidSummaryPMF|\\
\path|fno_oracleRecovery_expectedValidSummaryPMF|
\end{tabular}
\end{center}
These state that the true-oracle expected distortion is bounded by the
local-law budget measured through \(\widehat f\), plus the two-sided
proxy-oracle gap slack \(2\varepsilon_{\mathrm{orc}}\).

\subsection{Expanded DSL/IPW Certificate}

The translation table below gives the reading used by the main text. The
local-law argument supplies the validity, transport, oracle-transfer, and
certificate arithmetic. The survey-sampling layer supplies the standard
HT/H\'ajek/IPW, standard-error, law-of-large-numbers, and effective-sample-size
interpretation cited in the main text.

\begin{center}
\small
\setlength{\tabcolsep}{3pt}
\begin{tabular}{@{}p{0.18\linewidth}p{0.24\linewidth}p{0.29\linewidth}p{0.20\linewidth}@{}}
\toprule
Reader object & Econometric object & Audit notation & Paper role \\
\midrule
Local-law population &
Finite population of realized C1/C3/C2 checks &
\texttt{samples}, \(N,M,R\) in \texttt{computeDSLBound} &
Defines what the audit estimates. \\
\addlinespace
Audited law error &
HT/H\'ajek/IPW estimate of aggregate local-law residual &
\texttt{ipwUnionBound} &
Point estimate \(\widehat E_{\mathrm{law},R}^{\mathrm{IPW}}\). \\
\addlinespace
Audit uncertainty &
Design-based standard error for the same logged node sample &
\texttt{ipwUnionBoundSE} &
Sampling term \(\widehat{\mathrm{SE}}_{\mathrm{IPW}}\). \\
\addlinespace
Practical audit bound &
Estimate plus confidence multiplier and proxy-oracle gap margin &
\texttt{DSLBound.upperBound} &
Paper term
\(\widehat E_{\mathrm{law},R}^{\mathrm{IPW}}
+B_{\mathrm{gap}}+c_\alpha\widehat{\mathrm{SE}}_{\mathrm{IPW}}\). \\
\addlinespace
True-oracle transfer &
Uniform recovery of \(\fstar\) by \(\widehat f\) &
\texttt{OracleRecoveredWithin} &
Adds \(2\varepsilon_{\mathrm{orc}}\). \\
\addlinespace
Certified tree error &
Audit bound plus oracle-transfer slack &
\texttt{CertifiedTreeError} &
Full certificate. \\
\addlinespace
Expected-valid summary &
Aggregate approximate local laws imply expected root validity &
\path|approxLocalLaws_expectedValidSummaryPMF| &
Does not claim supportwise validity from aggregate audit bounds. \\
\addlinespace
Exact valid summaries &
Exact local laws imply root support lies in valid summaries &
\path|exactLocalLaws_root_validSummaryPMF|;\newline
\path|exactLocalLaws_multiround_validSummaryPMF| &
Exact supportwise validity. \\
\bottomrule
\end{tabular}
\end{center}

For sampled audits, we use the existing DSL/IPW certificate quantity rather than
introducing a new black-box IPW term. The canonical quantity is
\texttt{DSLBound}. Its upper bound expands as
\[
  \texttt{b.upperBound}
  =
  \texttt{b.gap\_estimate}
  + \texttt{b.bias\_margin}
  + \texttt{b.z\_score}\cdot\texttt{b.se}.
\]
For a tree audit this corresponds to
\[
  \widehat E_{\mathrm{law},R}^{\mathrm{IPW}}
  + B_{\mathrm{gap}}
  + c_\alpha\,\widehat{\mathrm{SE}}_{\mathrm{IPW}}.
\]
Here \(c_\alpha\) is only the paper notation for the confidence multiplier
stored as \texttt{b.z\_score}; we avoid the letter \(z\) because it is already
used for internal summary states.
The estimate and standard error come from the DSL/IPW layer after the local-law
checks have been chosen as the finite-population outcomes, following the
standard design-based estimators cited in the main text. Concretely,
\[
\begin{aligned}
  b &= \texttt{computeDSLBound}(\texttt{samples},N,M,R,\texttt{cal},c_\alpha),
  \\
  b.\texttt{gap\_estimate} &= \texttt{ipwUnionBound}(\texttt{samples},N,M,R),
  \\
  b.\texttt{se} &= \texttt{ipwUnionBoundSE}(\texttt{samples},N,M,R).
\end{aligned}
\]
The first line is the IPW/H\'ajek estimate of the union-bound composition
\(N\widehat p_{\mathrm{leaf}}+M\widehat p_{\mathrm{merge}}
+(R-1)\widehat p_{\mathrm{idemp}}\).  The second line is the corresponding
design-based standard error computed from the same logged samples. Thus
\(\widehat E_{\mathrm{law},R}^{\mathrm{IPW}}\) and
\(\widehat{\mathrm{SE}}_{\mathrm{IPW}}\) are not additional assumptions; they
are the estimate and uncertainty fields of the existing DSL certificate.
The true-oracle certificate adds only the already-defined proxy-oracle gap slack:
\[
  \texttt{CertifiedTreeError b }\varepsilon_{\mathrm{orc}}
  =
  \texttt{b.upperBound}
  + 2\varepsilon_{\mathrm{orc}}.
\]
The expanded equality is
\[
\texttt{CertifiedTreeError\_eq\_expanded}.
\]
Thus the certificate has no hidden IPW remainder; it is exactly
the reported estimate plus the proxy-oracle gap margin and the existing DSL/IPW
standard-error term.

The appendix also records the standard empirical-Bernstein alternative
\citep{MaurerPontil2009}. The radius
\[
  \texttt{UnionBoundEBRadius}
  =
  N r_{\mathrm{leaf}} + M r_{\mathrm{merge}} + (R-1)r_{\mathrm{idemp}}
\]
is defined by applying \texttt{empiricalBernsteinRadius} separately to the
weighted leaf, merge, and re-summary samples.  The variant
\[
\texttt{ipwUnionBound\_empirical\_bernstein\_from\_components\_expanded}
\]
states that componentwise empirical-Bernstein events imply the corresponding
union-bound error event with exactly this displayed radius.  Effective sample
size enters through these standard errors or radii; it is a statement about
audit uncertainty, not a guarantee that more training data monotonically lowers
the learned residual.

\subsection{Schedule Invariance}

The same-tree containment is the minimal correspondence. Classical mergeable
summaries often prove more: a valid root is obtained for every allowed merge
schedule. For ordered data, "allowed" means rebracketings that preserve leaf
order. For unordered sketches, additional commutativity can make permutations
valid too.

C-TreePO has the corresponding query-level bridge. If C1 and C3 hold on two
trees \(T\) and \(T'\) with the same leaves, then both roots have zero
expected oracle distortion, hence the schedules agree at the task level:
\[
\texttt{localLaw\_schedule\_bridge}.
\]
This is deliberately weaker than byte-identical state equality. Classical
associative sketches may give identical root states; C-TreePO gives the
equivalence needed by the oracle/query the task cares about. When
intermediate summaries are reused in a two-level plan, C2 supplies the stable
re-entry condition:
\[
\texttt{localLaw\_foldOfFolds\_bridge}.
\]

For ordered free-monoid summaries, the classical side has the stronger
Gibbons-style rebracketing result:
\[
\texttt{orderedClassicalSchedule\_bridge}.
\]
It gives identical state evaluation for two merge trees with the same ordered
stream data, without assuming commutativity. This is the state-equality
counterpart to C-TreePO's oracle-level schedule invariance.

\subsection{One-Way and Restricted Summaries}

One-way summaries are not automatically arbitrary-tree summaries. Their merge
operation absorbs a fresh suffix into an existing state, so the certified
topology is a sequential fold unless a separate full mergeability guarantee is
available. In the local-law language, such a method can be used on the
topologies its validity statement covers; it should not be cited as a full
binary-tree nesting guarantee. This is why GK-style one-way summaries are marked
sequential-only, while full state-level sketches such as HLL, Count-Min, and
KLL-style mergeable variants can be used in arbitrary binary merge trees when
their corresponding witnesses are present.
